% MAS1004 Assignment 1: long report.
% Student-only decisions and Problem 7 material remain explicitly marked.
\documentclass[11pt]{article}
\usepackage{mas1004}

\reporttitle{Assignment 1: Long report}
\studentname{[STUDENT: add your name]}
\studentid{[STUDENT: add your student ID]}

\begin{document}
\makeheader

% ---------------------------------------------------------------------------
\section*{Problem 1. Categories and images}

\subsection*{Why these categories}
\fillin{[STUDENT REVIEW: explain why you chose basketball jump shot, basketball layup, and basketball dunk, and why you care about distinguishing them.]}

\subsection*{What I expected the model to use}
\fillin{[STUDENT REVIEW: before training, record which visual features you expected the model to use, such as body pose, ball position, hoop position, clothing, colour, texture, or background.]}

\subsection*{Images per category}
The collection step obtained 150 images for each category: basketball\_dunk, basketball\_jump\_shot, and basketball\_layup. After the documented cleaning, the saved count was 123 basketball\_dunk images, 61 basketball\_jump\_shot images, and 65 basketball\_layup images; the complete count output is included in Problem 4.

% ---------------------------------------------------------------------------
\section*{Problem 2. Preparing an image}

\picturefile{figures/problem2_comparison.png}
\picturefile{figures/prepare_resize_crop_check.png}

The saved comparisons show the original image beside the inverse-normalized display of \texttt{prepare\_image}. The crop check documents the required preparation: resize the shorter side to 256, then take a centered 224 by 224 crop. In the documented example, the crop removes material from the left and right sides of the 1920 by 1080 original; the report image shows the crop region and the resulting 224 by 224 display.

\fillin{[STUDENT REVIEW: identify an image in the comparison where the center crop cuts off part of a feature you care about, and explain what was cut off.]}

% ---------------------------------------------------------------------------
\section*{Problem 3. First training runs}

\begin{center}\footnotesize
\begin{tabular}{llllrrrr}
\toprule
tag & start & trained & epochs & lr & trainable parameters & train acc. & test acc. \\
\midrule
run & ImageNet & all layers & 10 & 0.0001 & 11,178,051 & 99.4\% & 78.9\% \\
scratch & random & all layers & 10 & 0.0001 & 11,178,051 & 88.3\% & 54.4\% \\
frozen & ImageNet & last layer & 10 & 0.001 & 1,539 & 81.9\% & 58.9\% \\
five\_epochs & ImageNet & all layers & 5 & 0.0001 & 11,178,051 & 99.4\% & 73.3\% \\
\bottomrule
\end{tabular}
\end{center}

\picturefile{figures/run_curves.png}
\fillin{[STUDENT REVIEW: describe whether the run loss went down, whether it continued downward or flattened, and what the curves suggest.]}

The run finished with 99.4\% training accuracy and 78.9\% test accuracy, so the training accuracy was higher than the held-out test accuracy. \fillin{[STUDENT REVIEW: explain what you think this train--test gap means.]}

\begin{center}\picturefile[0.5\linewidth]{figures/run_confusion.png}\end{center}

The run confusion matrix has its largest two-way mix-up between basketball\_dunk and basketball\_layup: 5 dunk images were predicted as layup and 5 layup images were predicted as dunk. \fillin{[STUDENT REVIEW: say whether this surprised you and why.]}

The ImageNet-started all-layer run reached 78.9\% test accuracy, while the random-start all-layer run reached 54.4\%. The ImageNet-started frozen-last-layer run reached 58.9\%. \fillin{[STUDENT REVIEW: explain why you think the ImageNet and random results differ, and how close last-layer-only training is to all-layer training.]}

\subsection*{Saved run figures}
\picturefile{figures/scratch_curves.png}
\picturefile{figures/scratch_confusion.png}
\picturefile{figures/frozen_curves.png}
\picturefile{figures/frozen_confusion.png}
\picturefile{figures/five_epochs_curves.png}
\picturefile{figures/five_epochs_confusion.png}

\subsection*{Output of the runs}
\outputfile{outputs/run_output.txt}
\outputfile{outputs/scratch_output.txt}
\outputfile{outputs/frozen_output.txt}
\outputfile{outputs/five_epochs_output.txt}

% ---------------------------------------------------------------------------
\section*{Problem 4. Cleaning}

\subsection*{My rule}
\fillin{[STUDENT REVIEW: insert the cleaning rule exactly as you wrote it before cleaning.]}

\subsection*{Three kinds of junk}
The saved removal log and count output include, among other reasons, \texttt{Video thumbnail}, \texttt{Drawing}, and \texttt{Tutorial}. The saved near-copy sheet also visibly documents duplicate pairs and examples of non-photographic material. \fillin{[STUDENT REVIEW: select one picture for each of your three junk types and explain why each violates your rule.]}
\picturefile{figures/copies.png}
\picturefile{figures/basketball_dunk.png}
\picturefile{figures/basketball_jump_shot.png}
\picturefile{figures/basketball_layup.png}

\subsection*{What I removed}
\outputfile{outputs/clean_count.txt}

The saved suspects sheet contained 33 near-duplicate pairs. \fillin{[STUDENT REVIEW: state which image from each pair you removed or kept, using the exact pair list if your instructor requires it.]}

Among the first 20 ranked suspects in each class, the removal log shows 6 removed for basketball\_dunk, 16 removed for basketball\_jump\_shot, and 13 removed for basketball\_layup. \fillin{[STUDENT REVIEW: say whether these suspects included junk you had missed and whether you found junk not listed by the suspects tool.]}

\subsection*{Before and after cleaning}
\outputfile{outputs/compare.txt}

The saved comparison reports the original run at 78.9\% downloaded-test accuracy and 55.6\% new-image accuracy, and the saved clean run at 72.0\% downloaded-test accuracy and 50.0\% new-image accuracy. \fillin{[STUDENT REVIEW: explain why either accuracy may have gone down, remembering that cleaning changed the downloaded test set.]}

\picturefile{figures/clean_curves.png}
\begin{center}\picturefile[0.5\linewidth]{figures/clean_confusion.png}\end{center}

% ---------------------------------------------------------------------------
\section*{Problem 5. Images from a new source}

\fillin{[STUDENT REVIEW: identify where the new images came from and why you are sure they were not among the downloaded images.]}

\outputfile{outputs/overlap.txt}

The overlap output says that none of the 18 new images was a near copy of the 448 downloaded images; the closest similarity was 0.87, below the reported near-copy threshold of 0.95.

The saved comparison gives the downloaded-test and new-image accuracies side by side as follows: original \texttt{run}: 78.9\% and 55.6\%; clean: 72.0\% and 50.0\%. The later exported-model check records 12/18 = 66.7\% on new images; that complete output is in the appendix.

\subsection*{Five confident mistakes}
\picturefile{figures/problem5_preparation_comparison.png}

The five mistakes listed by the saved check output are:

\begin{itemize}
\item \texttt{IMG\_6103.jpg}: true basketball\_dunk; predicted basketball\_jump\_shot (94\%).
\item \texttt{IMG\_6084.jpg}: true basketball\_dunk; predicted basketball\_layup (89\%).
\item \texttt{IMG\_6077.jpg}: true basketball\_dunk; predicted basketball\_jump\_shot (70\%).
\item \texttt{IMG\_6108.jpg}: true basketball\_dunk; predicted basketball\_jump\_shot (69\%).
\item \texttt{IMG\_6063.jpg}: true basketball\_layup; predicted basketball\_jump\_shot (57\%).
\end{itemize}

\fillin{[STUDENT REVIEW: for each of the five pictures, explain what visual detail you think pushed the model toward its wrong class. Do not replace this with an unverified claim.]}

\fillin{[STUDENT REVIEW: compare the new images with the downloaded images in background, lighting, angle, distance, and other content in the frame.]}

% ---------------------------------------------------------------------------
\section*{Problem 6. The demo}

Address: \fillin{[STUDENT REVIEW: insert the published GitHub Pages URL. The saved check output says docs/model.onnx, docs/model.json, and docs/selftest.json were missing or not committed, so no verified published URL is available in the project files.]}

\fillin{[STUDENT REVIEW: record what happens when you show the demo an image that is none of the three categories.]}

% ---------------------------------------------------------------------------
\section*{Problem 7. What the model looks at}

\fillin{[STUDENT REVIEW: add one subsection for every hypothesis you chose. The required material is listed below; no hypotheses, experimental reasoning, exact prompts, or before/after model answers were supplied in the verified project outputs.]}

\subsection*{Guess 1 --- student material required}
\fillin{[STUDENT: hypothesis and evidence from your images or mistakes]}

\fillin{[STUDENT: exact change chosen and why it tests the hypothesis]}

My request to the agent, word for word:
\begin{quote}
\fillin{[STUDENT: paste the exact request]}
\end{quote}

\fillin{[STUDENT: add the original and changed images, and record the model answer before and after for every image]}

\fillin{[STUDENT: conclusion about the hypothesis and confidence]}

\subsection*{What I now think the model uses}
\fillin{[STUDENT REVIEW: summarize the conclusions from all Problem 7 experiments and state how sure you are.]}

% ---------------------------------------------------------------------------
\section*{Appendix. Output of check.py}
\outputfile{outputs/check.txt}

\end{document}
